\section{A Discontinuity Account of the Stability Gap}
\label{sec:account}

The previous section surveyed three findings that each touch a different face of the same object: the gap survives the joint-loss oracle (Section~\ref{subsec:stability_gap}), gradual transitions forget less (Section~\ref{subsec:blurry}), and a low-loss corridor demonstrably connects the single-task solution to the joint optimum (Section~\ref{subsec:lmc}). None of that prior work names a single mechanism that ties them together. This section sets out the account this thesis defends: the stability gap is, in its dominant components, the downstream symptom of the discontinuity at the task boundary introduced in Section~\ref{sec:introduction}. We make the claim precise by decomposing the gap into three additive contributors (Section~\ref{subsec:decomposition}) and then by contrasting the two ways one can intervene on the discontinuity (Section~\ref{subsec:pathfinding_vs_shaping}).

\subsection{Decomposing the Stability Gap}
\label{subsec:decomposition}

Previous work has identified various causes for the stability gap: the gradient imbalance (Section~\ref{subsec:stability_gap}), the finite-buffer estimator (Section~\ref{subsec:experience_replay}), and the trajectory effect (Section~\ref{subsec:stability_gap}). Here, we unify these into a single decomposition, breaking the gap down into three distinct forces that act at different stages of training. We will isolate and measure each empirically in Section~\ref{sec:methods}.

Consider the first step of vanilla Experience Replay (ER) at the start of a new task, $T_1$. Because the optimal parameters for the previous task ($\theta_0^*$) form a stationary point for the replay loss, the true replay gradient at this point is near zero. Consequently, the first parameter update is driven almost entirely by the new task:
\begin{equation}
    \theta_1 \approx \theta_0^* - \eta\left(\varepsilon(\theta_0^*) + g_{\text{new}}(\theta_0^*)\right),
    \label{eq:first_step}
\end{equation}
where $g_{\text{new}} := \nabla L_{\text{current}}$ and $\varepsilon$ is the zero-mean error of the finite-buffer estimator. This initial step triggers three distinct phenomena, each contributing to the subsequent performance dip on $T_0$:

\paragraph{Magnitude asymmetry ($G_{\text{mag}}$).} At $\theta_0^*$ the new-task gradient $g_{\text{new}}$ does not vanish --- the new task is still unlearned --- whereas the buffer sampling noise $\varepsilon$ is comparatively small and shrinks as the buffer grows. The first step is therefore dominated by the new task, $-\eta\, g_{\text{new}}$, and carries the iterate out of the replay basin. This is the same mechanism as the trajectory dip of Section~\ref{subsec:stability_gap}: any displacement away from $\theta_0^*$ raises the replay loss, so the dip appears at once --- only here it is driven by the unbalanced first step rather than an exact joint step. Because the imbalance is fixed by the relative gradient magnitudes at step one, this is the dominant contributor to the gap, and corresponds to the gradient-imbalance effect observed by \citet{delange2023continualevaluationlifelonglearning}.

\paragraph{Estimator variance ($G_{\text{est}}$).} In any individual training run, $\varepsilon \neq 0$, which adds a stochastic kick to $\theta_1$. As discussed in Section~\ref{subsec:experience_replay}, the replay buffer provides only a finite-sample estimate of the past data \cite{aljundi2019gradientbasedsampleselection}. This per-step sampling noise creates a baseline loss degradation that persists regardless of the learning schedule, as it is independent of how the loss terms are composed.

\paragraph{Trajectory effect ($G_{\text{traj}}$).} Even in an idealized scenario where magnitude asymmetry is corrected and buffer noise is eliminated ($\varepsilon \equiv 0$), a geometric penalty remains. As outlined in Section~\ref{subsec:stability_gap}, an instantaneous switch in objectives forces the iterate off the optimal path of the joint loss. Along this new trajectory, the replay loss is strictly higher than it would be at the joint-loss optimum $\theta_{\text{joint}}^*$. This is the structural contributor addressed by the envelope-theorem argument in Section~\ref{sec:theory}.

Crucially, these three components dominate at different timescales: $G_{\text{mag}}$ drives the massive disruption at step one, $G_{\text{traj}}$ dictates the curve over the first equilibration window, and $G_{\text{est}}$ provides a constant hum of noise across all steps.

\subsection{Path-Finding versus Landscape-Shaping}
\label{subsec:pathfinding_vs_shaping}
The decomposition suggests two ways to intervene, distinguished by what they do with the discontinuity rather than by what curvature information they use. \emph{Path-finding} accepts the discontinuity. The objective still jumps from $L_{\text{replay}}$ to $L_{\text{replay}} + L_{\text{current}}$ at the boundary, and the intervention reshapes \emph{how} the optimiser moves, replacing plain SGD with a preconditioned step $\theta_{t+1} = \theta_t - \eta\, M(\theta_t)\,\hat{g}_t$. Here $M$ is built from past-task information --- the inverse Fisher of EWC \cite{Kirkpatrick_2017}, the precision matrix of NCL \cite{kao2021naturalcontinuallearningsuccess}, or the feasible-set projection of GEM \cite{lopezpaz2022gradientepisodicmemorycontinual}. The path-finding methods of Section~\ref{subsec:path_finding} all sit here.

\emph{Landscape-shaping} removes the discontinuity. It modifies the objective in a neighbourhood of the boundary so the surface the optimiser descends is continuous through the transition, leaving the geometry of the step to SGD's natural step-size adaptation. The $\lambda$-curriculum is the canonical example: the gradient followed is $\hat{g}_t = \nabla\!\left(L_{\text{replay}} + \lambda(t)\, L_{\text{current}}\right)(\theta_t)$ with $\lambda(t)$ a continuous schedule from $0$ to $1$. Two consequences distinguish it from path-finding. A preconditioner that normalises the gradient breaks SGD's step-size adaptation, whereas $\lambda(t)$ multiplies the loss and flows through the backward pass to every parameter without rescaling the step; and the preconditioned step needs a per-step matrix--vector product, whereas the curriculum adds one scalar multiply. Section~\ref{sec:theory} shows that this continuity is enough to remove the trajectory contributor $G_{\text{traj}}$, and that it simultaneously tempers $G_{\text{mag}}$ by introducing $g_{\text{new}}$ gradually rather than all at once.
